\par\noindent\begin{minipage}{\linewidth}
\begin{summarybox}
  \textbf{\textcolor{CSummary}{一句话说明它的作用}}

  \textbf{三边确定，面积就确定。}海伦公式让你在\textbf{不必求角、不必求高}的情况下，直接从三条边算面积。

  \medskip
  \textbf{\textcolor{CSummary}{使用条件}}
  $a,b,c>0$，并且 $a+b>c$、$b+c>a$、$c+a>b$。

  \textbf{\textcolor{CSummary}{核心公式}}
  \[
   \boxed{p=\frac{a+b+c}{2},\qquad S=\sqrt{p(p-a)(p-b)(p-c)}}.
  \]
  \textbf{进一步可得：}
  \[
    r=\frac Sp,\qquad h_a=\frac{2S}{a},\qquad h_b=\frac{2S}{b},\qquad h_c=\frac{2S}{c}.
  \]
  此处 $h_a$ 表示以边 $a$ 为底时对应的高，其余同理。

  \textbf{\textcolor{CSummary}{为什么公式里出现四个因式？}}

  证明先得 $16S^2=(a+b+c)(a+b-c)(a-b+c)(-a+b+c)$。
  这四项来自\textbf{两次平方差分解}；再把它们各写成 $2p$ 或 $2(p-\text{一边})$，就得到海伦公式。

  \medskip
  \textbf{\textcolor{CSummary}{考场上怎样选方法？}}
  \begin{enumerate}[leftmargin=*,itemsep=0.3em]
   \item \textbf{只给三边：}先验三角形不等式，再算半周长，直接代入海伦公式。
   \item \textbf{已知两边及夹角：}优先用 $S=\tfrac12 ab\sin C$。
   \item \textbf{有现成的直角或很容易作出高：}优先试底乘高除以二。
   \item \textbf{已知面积、周长并有未知边：}反用海伦公式列方程；最后检查三边条件。
  \end{enumerate}

  \textbf{\textcolor{CSummary}{两种特殊情形帮你记住公式}}

  等边三角形：$a=b=c=\ell$，代入得 $S=\sqrt3\ell^2/4$。
  直角三角形：$a,b$ 为直角边、$c$ 为斜边，代入并用 $c^2=a^2+b^2$，得 $S=ab/2$。
  它们都包含在海伦公式中；算题时仍选择更省步骤的方法。

  \medskip
  \textbf{\textcolor{CSummary}{证明路径}}
  \textbf{入门：}勾股定理消去高 $\rightarrow$ 平方差 $\rightarrow$ 半周长。\\
  \textbf{步骤少：}内切圆分段 $x,y,z$ $\rightarrow$ 正切两角和 $\rightarrow$ $S=pr$。\\
  \textbf{常规：}正弦面积公式 $\rightarrow$ 余弦定理 $\rightarrow$ 平方差。

  \medskip
  \textbf{\textcolor{CSummary}{简单核对例子}}
  三边 $13,14,15$，半周长 $p=21$，面积
  $S=\sqrt{21\cdot8\cdot7\cdot6}=84$，内切圆半径 $r=84/21=4$。
\end{summarybox}
\end{minipage}\par
